\section{Fringing, sky subtraction and extraction}
\label{sec:fringe}

\subsection{What \banzai{}~1.0.1 does}
\label{sec:fringe_banzai}
The \texttt{FRINGE} extension of each \texttt{-2d} product is the map $\Phi(x,y)$ the frame was divided by. We find
that:
\begin{enumerate}
\item $\Phi$ is a whole normalised lamp flat (taken 2026 March 14/15, eight months before the science), not only its
fringe component: besides the fringes it contains the lamp's illumination profile along the slit (bands of
$\pm5$ per cent at 5000--7000\,\AA) and flat-field structure, and it is applied from $\approx5000$\,\AA{}
(Fig.~\ref{fig:master}).
\item The map is shifted \emph{along the slit} by a fitted offset (\texttt{L1FRNGOF}). Cross-correlating the maps
of two frames gives a shift equal to the difference of their offsets (17 rows for offsets $+9.0$ and $-8.5$). The
fitted offsets range from $-575$ to $+182$\,pixels (clustered near $\pm9$ and $\pm14$, aliases of the fringe
pattern), whereas flexure is 1--2\,pixels.
\item Where the shifted map is undefined it holds quantisation noise ($\sim10^{-6}$) or values up to $\pm3\times10^4$.
In 13 of the 50 frames (all with offsets between $-8.5$ and $-11.1$) the invalid edge falls inside the illuminated
order at 5200--5700\,\AA{}, the pixels there are divided by $\gg1$, and the 2-d sky model oscillates into stripes
that are extracted as 25--127 spikes (Fig.~\ref{fig:edge}).
\item Every frame is divided by the lamp's slit profile at the rows of the object, a frame-dependent and
wavelength-dependent factor of up to $\pm10$ per cent (15 per cent excess at 5500--7000\,\AA{} in a bright-moon frame).
\end{enumerate}

\subsection{Undoing the division and the master pattern}
\label{sec:fringe_master}
The frame before the division is recovered as
\begin{equation}
R(x,y)=\begin{cases} I(x,y)\,\Phi(x,y) & 0.5<\Phi<2\\ I(x,y) & |\Phi|<10^{-3}\ \text{(not applied)}\\
\text{masked} & \text{otherwise,}\end{cases}
\label{eq:undo}
\end{equation}
with $I$ the \texttt{SCI} image, and likewise for the uncertainty. Each frame's map is the lamp pattern
$\Pi$ shifted along the slit, $\Phi_f(x,y)=\Pi(x,y-o_f)$, so $\Pi$ is recovered by shifting $\Phi_f$ by $-o_f$
(linear interpolation; pixels with interpolation weight $<0.99$ discarded). The 61 maps with $|o_f|<40$ give
patterns that agree with their median $\bar\Pi$ with $r=0.998$--1.000. The fringe component is
\begin{equation}
\delta(x,y) = w(\lambda)\left[\frac{\bar\Pi(x,y)}{\mathcal{M}_{41}[\bar\Pi](x,y)}-1\right],\qquad
w(\lambda)=\min\!\left[\max\!\left(\frac{\lambda-6500}{300},0\right),1\right],
\label{eq:delta}
\end{equation}
where $\mathcal{M}_{41}$ is a running median of 41 pixels along the dispersion. This removes the slit-illumination
bands and the smooth flat structure and keeps the fringes (period $\approx11$\,pixels along the dispersion,
10.9 per cent rms at 7000--9500\,\AA; Fig.~\ref{fig:master}).

\subsection{Fringe fit}
\label{sec:fringe_fit}
\textit{Shift.} To first order $R=R_0(1+a\,\delta)$ with $R_0$ smooth along the dispersion on fringe scales. With
$\mathcal{H}[\cdot]$ the residual from a 21-pixel running median along $x$, for every trial shift $(\Delta x,\Delta y)$
on a 0.5-pixel grid with $|\Delta x|,|\Delta y|\le2$,
\begin{equation}
\hat a=\frac{\sum w\,\mathcal{H}[R]\,\mathcal{H}[R\,\delta']}{\sum w\,\mathcal{H}[R\,\delta']^2},\qquad
\chi^2(\Delta x,\Delta y)=\sum w\left(\mathcal{H}[R]-\hat a\,\mathcal{H}[R\,\delta']\right)^2,
\label{eq:fringe2d}
\end{equation}
where $\delta'(x,y)=\delta(x-\Delta x,y-\Delta y)$, $w=1/\sigma_R^2$, and the sums run over the red-order pixels at
7000--9800\,\AA{} more than 6 rows from the slit ends. The minimum is well defined (period $\approx11$\,pixels in
$\Delta x$) and lies at $|\Delta x|\le1$, $|\Delta y|\le1$ for all frames.

\textit{Amplitude (first version).} The 2-d fit under-corrects, so the correction is applied to the extracted spectrum. The profile-weighted fringe of each 1-pixel bin $\beta$ is
\begin{equation}
f_\beta=\frac{\sum_{i\in\beta}P_i\,\delta'_i}{\sum_{i\in\beta}P_i},
\label{eq:fringe1d}
\end{equation}
and the extracted counts are divided by $1+A(\lambda)f_\beta$, where the shift of $\delta'$ is the 2-d one.
The amplitude $A(\lambda)$ can be fitted freely on a spectrum by linear least squares of
$\mathcal{H}_{\rm 1d}[F]$ on $\mathcal{H}_{\rm 1d}[F\,f]$ ($\mathcal{H}_{\rm 1d}$: difference of 25\,\AA{} and 200\,\AA{}
running medians, telluric bands excluded) in a wavelength region. This is reliable on the smooth standards but not on
a single AGN spectrum, where it latches onto real structure. The amplitude is, however, very stable within one
instrument configuration, so it is calibrated per slit width as the median of the free fits over an ensemble of
frames, in four regions (6800--7500, 7500--8300, 8300--9100, 9100--9800\,\AA), and interpolated linearly between their
centres (Table~\ref{tab:amplaw}). Regions where the median is negative or its scatter exceeds half its value
(6800--7500\,\AA, where $\delta$ is still tapered in and most pixels are telluric) take the value of the nearest
reliable region. This uses no target information. The amplitude exceeds unity because the stored lamp map has a
diluted fringe contrast, and it differs between the 2 and 6\,arcsec slits, plausibly because slit-filling light reaches
each pixel over a wider wavelength range in a wide slit. This model is kept as an option
(\texttt{method='amplitude'}); Section~\ref{sec:fringe_harmonic} explains why it leaves $\approx5$ per cent
residuals and describes the default model.

\begin{table}
\centering
\caption{Fringe amplitude $A(\lambda)$ per slit width: median and robust scatter of the free fits over the ensemble of frames. $^{a}$Unreliable region (negative or scatter $>$ half the value); the value of the nearest reliable region is used.}
\label{tab:amplaw}
\begin{tabular}{lcccc}
\toprule
Slit & 7150\,\AA & 7900\,\AA & 8700\,\AA & 9450\,\AA \\
\midrule
6\,arcsec (50 frames) & \AmpSixA & \AmpSixB & \AmpSixC & \AmpSixD \\
2\,arcsec (18 frames) & \AmpTwoA & \AmpTwoB & \AmpTwoC & \AmpTwoD \\
\bottomrule
\end{tabular}
\end{table}


\subsection{Harmonic fringe model}
\label{sec:fringe_harmonic}
\textit{Why a scaled lamp pattern cannot work.} Fig.~\ref{fig:harmonic} shows the red order of a standard and of a
Fairall~9 frame before fringe correction, divided by a smooth continuum: the fringes have sharp maxima and broad,
shallow minima, as expected from the Fabry--P\'erot transmission of the thinned CCD. The lamp pattern scaled by
$A(\lambda)$ is almost sinusoidal; matching the maxima makes its minima too deep (0.55--0.6 against 0.75--0.8), which
produces the asymmetric, spiky residuals of the amplitude-law correction. The reason is the spectral bandwidth that
reaches each pixel. A pixel at $x$ receives wavelengths $\lambda_0(x)-D\,s$ from every slit position $s$ that is
illuminated, so the recorded fringe is the transmission $\mathcal{T}(x,y,\lambda)$ smoothed by the image of the light
across the slit, in wavelength units: a top-hat of the slit width (51\,\AA{} in the red order for 6\,arcsec, larger than
the $\approx38$\,\AA{} fringe period) for the lamp, and the seeing profile for a point source. Writing
$\mathcal{T}$ locally as a Fourier series in its phase $\phi$, harmonic $k$ is multiplied by the Fourier transform of
the kernel at $k$ times the fringe frequency. A symmetric kernel preserves the phase, but the slit-filling lamp
suppresses the fundamental and, much more, the harmonics, by factors that depend on wavelength (the period grows as
$\lambda^2$) and differ from the object's. The lamp flat therefore gives the \emph{phase} of the fringes but neither
their amplitude nor their shape. The object's phase is also offset from the lamp's by flexure and by its position
across the wide slit, which changes the wavelength falling on each pixel: a 5\,\AA{} offset moves the phase by
$\approx0.2$ cycles. The free fits below give a frame-to-frame phase scatter of 16--26$^\circ$ for Fairall~9 and
6--9$^\circ$ for EG\,21; a 20$^\circ$ error alone leaves a residual of $\approx2\sin(10^\circ)A\approx10$ per cent of the
fringe amplitude.

\textit{Phase map.} For every row of the red order, the phase is the argument of the analytic signal of the fringe
component (lightly smoothed, Gaussian $\sigma=1$ pixel):
\begin{equation}
\phi(x,y)=\arg\left\{\tilde\delta(x,y)+i\,\mathcal{Q}_x[\tilde\delta](x,y)\right\},
\label{eq:phase}
\end{equation}
where $\mathcal{Q}_x$ is the Hilbert transform along the dispersion and $\tilde\delta$ the smoothed $\delta$ (eq.~\ref{eq:delta}).

\textit{Model and fit.} The fringe transmission of the extracted spectrum is modelled in logarithm as
\begin{equation}
\ln T(\lambda)=\sum_{k=1}^{K}\sum_{j=1}^{J}h_j(\lambda)\left[a_{jk}\,C_{k,\beta}+b_{jk}\,S_{k,\beta}\right],
\label{eq:harmonic}
\end{equation}
\begin{equation}
C_{k,\beta}=\frac{\sum_{i\in\beta}P_iw_i\cos k\phi_i}{\sum_{i\in\beta}P_i},
\label{eq:harmproj}
\end{equation}
with $S_{k,\beta}$ defined likewise with $\sin k\phi_i$, $w_i$ the taper of eq.~(\ref{eq:delta}), and $h_j$ linear
interpolation functions on the nodes $\lambda_j=6800,7300,\dots,9800$\,\AA{} ($J=7$). We use $K=2$ (28 coefficients);
$K=1$ is only slightly worse. The $b_{jk}$ absorb the phase offset of the frame and the $a_{jk},b_{jk}$ together give
the amplitude and shape, all smooth in wavelength. The coefficients are the linear least-squares solution of
\begin{equation}
\mathcal{H}_{150}[\ln F_\beta]=\sum_{j,k}a_{jk}\,\mathcal{H}_{150}[h_jC_k]_\beta+b_{jk}\,\mathcal{H}_{150}[h_jS_k]_\beta,
\label{eq:harmfit}
\end{equation}
where $F$ is the extraction without fringe correction and $\mathcal{H}_{150}$ subtracts a 150\,\AA{} running median.
Three iterations of $5\sigma$ clipping are applied. The telluric bands need care: if they are simply excluded,
the coefficients across 8900--9300\,\AA{} are constrained only by the neighbouring regions and the residual there
is 8--9.5 per cent. The fit is therefore done in two passes: the first excludes
the bands (6860--6960, 7150--7350, 7560--7720, 8100--8380 and 8900--9300\,\AA) and is followed by the wavelength and
flux calibration, which yield the telluric transmission $t$ of each spectrum (Section~\ref{sec:flux_tell}); the second
refits eq.~(\ref{eq:harmfit}) on $F/t$, excluding only the A band and pixels with $t<0.3$. In a test on all frames this
lowers the residual at 8950--9300\,\AA{} from 8--9.5 to 5.3--5.6 per cent and at 9300--9700\,\AA{} from 8.2--8.5 to
6.8--8.5 per cent, with the clean windows unchanged. The extracted counts and their
uncertainties (optimal and fixed aperture) are divided by $T$. The fit uses only the frame's own spectrum and the phase
map, and no information on the target. Because the basis oscillates with the fringe phase, it is nearly orthogonal to
astrophysical structure on scales above $\approx50$\,\AA, and the clipping removes narrow features.

\textit{Variants tested.} (i) The phase can be taken from the eight-month-old master pattern or from the lamp flat
taken in the same block as each science frame (69 such flats exist for the Fairall~9 campaign, matched by block
identifier; option \texttt{--flat-dir}). Their fringe components correlate at 0.95--0.98 with the master, and the
harmonic correction with same-block phases is indistinguishable from the one with the master (Table~\ref{tab:fringevariants}).
The detector pattern is stable over the eight months, so the old flat was not the limiting factor. (ii) The classical
approach of measuring the fringes on the standard and transferring them to the science (here: the high-passed
$\ln F$ of the nearest EG\,21, with a free shift and a free scale at each node) does better than the amplitude law
but worse than the harmonic model, because the phase offset differs between exposures. (iii) Fixing the harmonic
coefficients to their ensemble median, or fixing only the amplitude and fitting the phase, is worse than free per-frame
coefficients for the AGN.

\begin{table}
\centering
\caption{Fringe residuals of the red order (rms of the high-passed $\ln F$ in 7750--8050 and 8450--8900\,\AA, per cent;
median over frames; no short smoothing, so photon noise contributes the floor given). $^{a}$16 frames with a lamp flat
within one day.}
\label{tab:fringevariants}
\begin{tabular}{lcc}
\toprule
Correction & Fairall~9 & EG\,21 \\
\midrule
None & \FVnoneF & \FVnoneE \\
Amplitude law (master pattern) & \FVampF & \FVampE \\
Standard-star transfer & \FVstdF & -- \\
Harmonic, ensemble coefficients & \FVensF & \FVensE \\
Harmonic, master phase & \FVharmF & \FVharmE \\
Harmonic, same-block flat phase & \FVflatF & \FVflatE$^{a}$ \\
Photon-noise floor & \FVnoiseF & \FVnoiseE \\
\bottomrule
\end{tabular}
\end{table}

\subsection{Sky model}
\label{sec:sky}
For every 1-pixel wavelength bin $\beta$ of an order (\banzai{}'s binning, which follows the line tilt), the sky is a
polynomial along the slit,
\begin{equation}
S_\beta(y_{\rm o})=\sum_{k=0}^{2}b_{\beta k}\,y_{\rm o}^k,
\label{eq:sky}
\end{equation}
fitted by weighted least squares to the unmasked pixels with $|y_{\rm p}|>m\,\sigma_{\rm p}$ and
$|y_{\rm o}|<y_{\rm max}-6$, with iterative $3\sigma$ clipping. The mask half-width is $m=4$ unless that leaves fewer
than 25 per cent of the slit's pixels, in which case $m$ is reduced to the corresponding quantile of
$|y_{\rm p}|/\sigma_{\rm p}$ (but not below 2); the degree is reduced while fewer than $4(K+1)$ sky pixels remain, and set
to zero when all sky pixels lie on one side of the object. When the mask has been reduced, the sky pixels contain the
object's wings, so the sky polynomial and the object are fitted jointly to all slit pixels,
$R_i=\sum_k b_{\beta k}y_{{\rm o},i}^k+F'_\beta P_i$, with iterative clipping; only the sky part is kept. In synthetic
tests with $\sigma_{\rm p}=9$\,pixels in a 60-row slit this removes an 11 per cent flux bias. Its covariance
$C_\beta=(V^{\rm T}WV)^{-1}$ ($V$ the Vandermonde matrix, $W$ the inverse variances) is propagated into the
extraction errors. This replaces the 2-d spline sky of \banzai{} \citep{Kelson2003}, which cannot follow the
discontinuities introduced by the fringe map.

\subsection{Extraction}
\label{sec:extraction}
For bin $\beta$, with net counts $d_i=R_i-S_\beta(y_{{\rm o},i})$, profile weights $P_i$ (\banzai's fitted profile,
normalised to unit sum per column) and pixels within $|y_{\rm p}|<5\sigma_{\rm p}$, the optimal extraction
\citep{Horne1986} is
\begin{equation}
F_\beta=\frac{\sum_i a_i d_i}{\sum_i a_i P_i},\qquad a_i=\frac{P_i}{V_i},\qquad
V_i = r^2+\max(F_\beta P_i+S_{\beta,i},0),
\label{eq:horne}
\end{equation}
iterated with $V$ from the model (not the data); at each iteration the single pixel with the largest
$(d_i-F_\beta P_i)^2/V_i$ is rejected if this exceeds 25 (cosmic rays), up to $\max(2,\,0.1N_\beta)$ pixels per bin
(more ``outliers'' indicate a profile mismatch, not cosmic rays). The read noise is $r=3.5\,e^-$. Its variance is
\begin{equation}
\sigma^2_{F_\beta}=\frac{\sum_i a_i^2V_i+\mathbf{g}^{\rm T}C_\beta\,\mathbf{g}}{\left(\sum_i a_iP_i\right)^2},\qquad
\mathbf{g}=\sum_i a_i\,\mathbf{v}_i,
\label{eq:hornevar}
\end{equation}
with $\mathbf{v}_i=(y_{{\rm o},i}^2,y_{{\rm o},i},1)$ the sky basis at pixel $i$. A fixed-aperture extraction for RM is
\begin{equation}
F^{\rm box}_\beta=\frac{\sum_{|y_{\rm p}|<h}d_i}{\sum_{\rm all\ rows}P_i},
\label{eq:boxcar}
\end{equation}
with $h=10$\,pixels ($\pm3.4$\,arcsec) by default; the denominator is the number of columns the bin spans, which
gives the same normalisation as \banzai{} for bins that collect 0, 1 or 2 pixels per row. In synthetic tests the
recovered fluxes are unbiased (optimal 2 per cent, boxcar 3 per cent including a cosmic ray on the trace) and the
normalised residuals have rms 1.03 (optimal) and 1.09 (boxcar).
